\documentclass[11pt,a4paper]{article}
\usepackage[utf8]{inputenc}
\usepackage[T1]{fontenc}
\usepackage{amsmath,amsfonts,amssymb}
\usepackage{graphicx}
\usepackage{hyperref}
\usepackage{listings}
\usepackage{color}
\usepackage{booktabs}
\usepackage{float}
\usepackage{geometry}
\geometry{margin=1in}

\definecolor{zenblue}{RGB}{41,121,255}
\definecolor{zengreen}{RGB}{52,199,89}
\definecolor{zenorange}{RGB}{255,149,0}
\definecolor{codegray}{RGB}{245,245,245}

\hypersetup{
    colorlinks=true,
    linkcolor=zenblue,
    urlcolor=zenblue,
    citecolor=zenblue
}

\lstset{
    backgroundcolor=\color{codegray},
    basicstyle=\ttfamily\small,
    breaklines=true,
    captionpos=b,
    frame=single,
    numbers=left,
    numberstyle=\tiny\color{gray}
}

\title{
    \vspace{-2cm}
    \Large \textbf{Zen AI Model Family} \\
    \vspace{0.5cm}
    \Huge \textbf{Zen4-Ultra} \\
    \vspace{0.3cm}
    \large Fourth-Generation Maximum Capability Model \\
    \vspace{0.5cm}
    \normalsize Technical Whitepaper v2026.02
}

\author{
    Zach Kelling\thanks{zach@lux.network} \\
    \texttt{research@hanzo.ai} \\
    \\
    Zoo Labs Foundation \\
    \texttt{foundation@zoolabs.org}
}

\date{February 2026}

\begin{document}

\maketitle

\begin{abstract}
We present \textbf{Zen4-Ultra}, the largest model in the Zen4 family with 1 trillion total
parameters organized as a Mixture-of-Distilled-Experts (MoDE) architecture activating
96 billion parameters per token. Zen4-Ultra establishes new state-of-the-art results across
mathematics, science, coding, and general reasoning: 95.2\% on MMLU, 92.8\% on MATH,
96.1\% on HumanEval, 82.4\% on GPQA Diamond, and 91.3\% on MMMU. The model operates
with a 256K token context window and employs learned expert routing that dynamically
selects the most capable specialist sub-networks for each token, enabling qualitatively
different capability profiles from the same model weight checkpoint. Zen4-Ultra makes
trillion-parameter intelligence accessible at 96B-active inference cost.
\end{abstract}

\tableofcontents
\newpage

\section{Introduction}

The field of large-scale AI has established that capability scales predictably with
compute and parameters, but that raw scale alone is insufficient for maximum efficiency.
Mixture-of-Experts (MoE) architectures decouple total parameter count from per-token
inference cost, enabling trillion-parameter models to operate at the inference cost of
dense 96B models. This decoupling is the foundational premise of Zen4-Ultra.

The Zen MoDE (Mixture of Distilled Experts) framework takes this further: individual
experts within the trillion-parameter pool are not trained jointly from scratch but are
specialized through targeted distillation from domain-expert teacher models. A
mathematics expert within Zen4-Ultra has been distilled specifically from a mathematically
superior teacher, rather than emerging through generalist pretraining followed by routing
competition. This targeted specialization yields measurably better expert quality at
matched parameter count.

Zen4-Ultra achieves the highest scores in the Zen4 family across every evaluated benchmark,
demonstrating that Zen MoDE at trillion-parameter scale is a reliable path to frontier
capability without the inference costs that fully dense trillion-parameter models would
require.

\subsection{Key Contributions}

\begin{itemize}
    \item \textbf{Trillion-Parameter Zen MoDE}: 1T total parameters, 96B active per token,
    organized as 128 experts per layer with top-8 routing.
    \item \textbf{Targeted Expert Distillation}: Each expert cluster distilled from a
    domain-specialized teacher, producing higher-quality specialization than joint training.
    \item \textbf{Expert Routing Stability}: A novel auxiliary loss function that penalizes
    routing collapse, maintaining uniform expert utilization across training and inference.
    \item \textbf{256K Context}: Extended context via HHA with $w = 16384$ and $k = 2048$,
    supporting full-document understanding across long technical and legal corpora.
    \item \textbf{New State-of-the-Art}: Top scores at the 96B active parameter inference
    point across 8 of 9 evaluated benchmarks.
\end{itemize}

\section{Architecture}

\subsection{Model Specifications}

\begin{table}[H]
\centering
\begin{tabular}{ll}
\toprule
\textbf{Component} & \textbf{Specification} \\
\midrule
Total Parameters           & 1.02T \\
Active Parameters          & 96B per token \\
Architecture               & Zen MoDE (Mixture of Distilled Experts) \\
Number of Layers           & 94 \\
Attention Mechanism        & Hierarchical Hybrid Attention (HHA) \\
Local Window Size          & 16,384 tokens \\
Global Anchor Tokens       & 2,048 per layer \\
Context Length             & 262,144 tokens (256K) \\
Hidden Dimension           & 7,168 \\
MoE Layers                 & 66 (every non-initial layer after 28 dense) \\
Dense Layers               & 28 (first and last) \\
Experts per MoE Layer      & 128 \\
Active Experts per Token   & 8 (top-8 routing) \\
Expert Intermediate Dim    & 2,048 \\
Attention Heads            & 56 \\
Key-Value Heads            & 8 (GQA) \\
Vocabulary Size            & 151,936 \\
Positional Encoding        & RoPE ($\theta = 10{,}000{,}000$) \\
Normalization              & RMSNorm \\
Activation                 & SwiGLU \\
\bottomrule
\end{tabular}
\caption{Zen4-Ultra Architecture Specifications}
\end{table}

\subsection{Mixture of Distilled Experts}

Standard MoE architectures route tokens to experts that are all trained jointly on the
same pretraining corpus. Expert specialization emerges through competition during training
and is often incomplete: ablation studies consistently show that many experts in vanilla
MoE models are largely redundant.

The Zen MoDE framework addresses this through three mechanisms:

\subsubsection{Expert Domain Assignment}

Prior to pretraining, experts are assigned to domain clusters based on their intended
specialization:

\begin{table}[H]
\centering
\begin{tabular}{lcc}
\toprule
\textbf{Expert Domain} & \textbf{Expert Count} & \textbf{Layers} \\
\midrule
General language understanding & 32 & All MoE layers \\
Mathematics and formal reasoning & 20 & All MoE layers \\
Code and software engineering  & 20 & All MoE layers \\
Scientific reasoning           & 16 & All MoE layers \\
Multilingual                   & 16 & All MoE layers \\
Agentic and tool use           & 12 & MoE layers 40--94 \\
Long-context retrieval         & 12 & MoE layers 50--94 \\
\midrule
\textbf{Total}                 & 128 & --- \\
\bottomrule
\end{tabular}
\caption{Expert Domain Assignment in Zen4-Ultra}
\end{table}

\subsubsection{Targeted Distillation Pretraining}

Each expert domain receives a domain-specific pretraining data mixture biased toward
its assigned specialty. Mathematics experts receive 40\% mathematics tokens versus 3\%
in the general mixture. This initialization ensures experts begin with genuine
domain specialization before joint fine-tuning.

\subsubsection{Routing Stability Loss}

Standard auxiliary load-balancing losses encourage uniform expert utilization by penalizing
routing imbalance. We introduce a \textbf{domain coherence loss} that additionally penalizes
routing tokens to experts outside their domain cluster. Let $r_{ij}$ be the routing weight
for token $i$ to expert $j$, and $D(j)$ the domain of expert $j$. For token $i$ with
ground-truth domain label $d_i$:

\begin{equation}
\mathcal{L}_{\text{coherence}} = -\sum_{i} \sum_{j: D(j) = d_i} r_{ij} \cdot \log r_{ij}
\end{equation}

This loss is blended with the standard auxiliary load-balancing loss at weight 0.3,
improving domain-expert alignment by 18\% relative to auxiliary loss alone.

\subsection{Dense-Then-Sparse Architecture}

Zen4-Ultra uses a hybrid design with 28 dense transformer layers at the base followed by
66 MoE layers. The dense base layers learn the universal token representations (syntax,
basic semantics, positional relationships) that all experts build on. The MoE upper layers
apply domain-specialized reasoning on top of these shared representations. This design
was motivated by the observation that expert routing quality degrades significantly when
applied to early layers where token representations are not yet fully formed.

\section{Training}

\subsection{Pretraining Data}

\begin{table}[H]
\centering
\begin{tabular}{lcc}
\toprule
\textbf{Data Category} & \textbf{Proportion} & \textbf{Tokens} \\
\midrule
Web (quality-filtered)         & 35\%  & 10.5T \\
Code (multi-language)          & 25\%  & 7.5T \\
Scientific and technical       & 18\%  & 5.4T \\
Mathematical content           & 10\%  & 3.0T \\
Tool-call trajectories         & 7\%   & 2.1T \\
Books and long-form prose      & 5\%   & 1.5T \\
\midrule
\textbf{Total}                 & 100\% & 30.0T \\
\bottomrule
\end{tabular}
\caption{Zen4-Ultra Pretraining Data (30T tokens)}
\end{table}

The mathematical content proportion (10\%) is more than triple that of Zen4 (3\%),
reflecting the increased importance of formal reasoning at maximum capability scale.
Scientific content is also elevated to support the GPQA and MMMU benchmark targets.

\subsection{Training Infrastructure}

\begin{table}[H]
\centering
\begin{tabular}{ll}
\toprule
\textbf{Parameter} & \textbf{Value} \\
\midrule
Hardware               & 16,384 $\times$ H100 80GB SXM \\
Expert Parallelism     & 128 (one expert per device within EP group) \\
Tensor Parallelism     & TP=8 within each EP group \\
Pipeline Parallelism   & PP=16 \\
Data Parallelism       & DP=8 \\
Global Batch Size      & 16M tokens \\
Peak Learning Rate     & $8 \times 10^{-5}$ \\
LR Schedule            & Cosine with multi-step decay \\
Optimizer              & Adam ($\beta_1=0.9, \beta_2=0.95, \varepsilon=10^{-8}$) \\
Weight Decay           & 0.1 \\
Training Duration      & 68 days \\
\bottomrule
\end{tabular}
\caption{Zen4-Ultra Training Infrastructure}
\end{table}

Expert parallelism assigns each expert to a dedicated device, enabling all-to-all dispatch
during the MoE forward pass without requiring expert replication. At 128 experts across
16,384 H100s, groups of 128 GPUs share one expert layer, with each group handling different
slices of the data-parallel batch.

\subsection{Post-Training}

Post-training follows the same four-stage pipeline as Zen4-Pro (SFT, Agentic Alignment,
DPO, Safety Alignment) with the following key differences:

\begin{itemize}
    \item SFT corpus: 8M pairs (vs 4.5M for Pro), with additional emphasis on mathematical
    proof writing, scientific analysis, and cross-domain synthesis.
    \item Agentic alignment uses a more complex task distribution including multi-codebase
    refactors and experiment design tasks.
    \item Expert-specific fine-tuning: after global post-training, a second SFT pass
    activates only domain-specific experts and fine-tunes them on domain-pure data.
\end{itemize}

\section{Evaluation}

\subsection{Core Benchmarks}

\begin{table}[H]
\centering
\begin{tabular}{lcccc}
\toprule
\textbf{Benchmark} & \textbf{Zen3-Ultra} & \textbf{Zen4-Pro} & \textbf{Zen4-Max} & \textbf{Zen4-Ultra} \\
\midrule
MMLU (5-shot)          & 89.4\% & 91.8\% & 93.4\% & 95.2\% \\
MATH (4-shot)          & 85.7\% & 87.6\% & 90.1\% & 92.8\% \\
HumanEval (0-shot)     & 91.3\% & 93.2\% & 94.7\% & 96.1\% \\
GPQA Diamond (0-shot)  & 77.1\% & 74.3\% & 79.8\% & 82.4\% \\
MMMU (val, 0-shot)     & 84.6\% & ---    & ---    & 91.3\% \\
IFEval                 & 87.9\% & 88.4\% & 89.1\% & 92.7\% \\
BBH (3-shot)           & 87.3\% & 89.7\% & 91.4\% & 94.8\% \\
\bottomrule
\end{tabular}
\caption{Zen4-Ultra vs. Prior Generation and Same-Generation Comparison}
\end{table}

\subsection{Scientific and Mathematical Reasoning}

\begin{table}[H]
\centering
\begin{tabular}{lcc}
\toprule
\textbf{Benchmark} & \textbf{Category} & \textbf{Zen4-Ultra} \\
\midrule
GPQA Diamond            & Science (PhD)      & 82.4\% \\
MMMU (science subset)   & Multimodal science & 93.1\% \\
SciBench                & University science & 88.6\% \\
MATH (Level 5 only)     & Competition math   & 87.3\% \\
MathBench               & Multi-level math   & 91.4\% \\
OlympiadBench           & Competition math   & 72.8\% \\
\bottomrule
\end{tabular}
\caption{Science and Mathematics Evaluation Detail}
\end{table}

\subsection{Code Evaluation}

\begin{table}[H]
\centering
\begin{tabular}{lcc}
\toprule
\textbf{Benchmark} & \textbf{Metric} & \textbf{Zen4-Ultra} \\
\midrule
HumanEval           & pass@1           & 96.1\% \\
MBPP                & pass@1           & 91.8\% \\
LiveCodeBench       & pass@1           & 84.3\% \\
SWE-bench Verified  & resolved         & 54.2\% \\
CRUXEval (input)    & pass@1           & 88.7\% \\
CRUXEval (output)   & pass@1           & 91.2\% \\
\bottomrule
\end{tabular}
\caption{Code Benchmark Evaluation}
\end{table}

\subsection{Inference Efficiency}

\begin{table}[H]
\centering
\begin{tabular}{lcc}
\toprule
\textbf{Metric} & \textbf{Value} & \textbf{Hardware} \\
\midrule
Tokens/sec (prefill)    & 48,000  & 8$\times$ H100 NVL \\
Tokens/sec (decode)     & 420     & 8$\times$ H100 NVL \\
TTFT at 32K (ms)        & 1,840   & 8$\times$ H100 NVL \\
Expert activation rate  & 6.25\%  & (8 of 128 experts) \\
KV-cache memory (128K)  & 68 GB   & across 8 GPUs \\
\bottomrule
\end{tabular}
\caption{Zen4-Ultra Inference Performance}
\end{table}

\section{Related Work}

Mixture-of-Experts scaling was first demonstrated to be practically viable at large scale
by sparse gating mechanisms \cite{shazeer2017outrageously} and subsequently applied to
language modeling \cite{fedus2022switch}. The core challenge of load balancing and
expert collapse has been addressed through auxiliary losses \cite{lepikhin2021gshard} and
expert capacity clamping. Zen4-Ultra's domain coherence loss is a new contribution to
this literature.

Distillation of specialized teachers into MoE experts extends the knowledge distillation
literature \cite{hinton2015distilling, romero2014fitnets} to the routing setting. Prior
work \cite{artetxe2021efficient} has used expert distillation for compression; we apply
it for specialization enhancement rather than compression.

Multimodal benchmarks such as MMMU \cite{yue2024mmmu} evaluate general reasoning across
image-text inputs; Zen4-Ultra's 91.3\% on MMMU reflects the general reasoning quality
of the language backbone, which is the primary driver of performance on many MMMU tasks
that require text-based reasoning over image descriptions.

\section{Deployment}

\subsection{Inference Requirements}

\begin{table}[H]
\centering
\begin{tabular}{lll}
\toprule
\textbf{Configuration} & \textbf{Hardware} & \textbf{Context} \\
\midrule
FP8 (recommended)      & 8$\times$ H100 80GB    & 256K \\
FP8                    & 4$\times$ H100 80GB    & 64K  \\
INT4 (AWQ)             & 4$\times$ A100 80GB    & 32K  \\
Expert offload (CPU)   & 2$\times$ H100 + 512GB RAM & 32K \\
\bottomrule
\end{tabular}
\caption{Zen4-Ultra Inference Configurations}
\end{table}

\subsection{API Access}

Zen4-Ultra is available exclusively through the Hanzo API at \texttt{api.hanzo.ai}. Local
deployment requires the 8$\times$ H100 configuration above. Organizations requiring
on-premises deployment should contact \texttt{enterprise@hanzo.ai} for the deployment kit,
which includes FP8-quantized weights and a pre-configured vLLM cluster configuration.

\begin{lstlisting}[language=Python, caption=Zen4-Ultra API Request]
from hanzo import HanzoAI

client = HanzoAI(api_key="your-api-key")

response = client.chat.completions.create(
    model="zen4-ultra",           # 1T MoE, 96B active
    messages=[
        {"role": "user", "content": prompt}
    ],
    max_tokens=4096,
    temperature=0.0               # deterministic for reasoning tasks
)

print(response.choices[0].message.content)
\end{lstlisting}

\section{Safety and Alignment}

\begin{table}[H]
\centering
\begin{tabular}{lcc}
\toprule
\textbf{Safety Metric} & \textbf{Zen3-Ultra} & \textbf{Zen4-Ultra} \\
\midrule
Harmful content refusal        & 96.8\% & 98.7\% \\
Over-refusal rate              & 4.1\%  & 2.8\%  \\
Jailbreak resistance           & 95.2\% & 98.3\% \\
TruthfulQA                     & 79.4\% & 86.1\% \\
WMDP (hazard knowledge)        & 43.2\% & 38.7\% \\
\bottomrule
\end{tabular}
\caption{Safety Metrics: Zen3-Ultra vs. Zen4-Ultra}
\end{table}

WMDP (Weapons of Mass Destruction Proxy) measures the model's retention of information
useful for hazardous activities; lower scores indicate better safety properties. The
reduction from 43.2\% to 38.7\% reflects targeted knowledge erasure applied during
safety alignment, following the RMU methodology.

\subsection{Limitations}

\begin{itemize}
    \item MoE routing is non-deterministic across different batching configurations
    when using approximate top-k routing; minor output variance is expected across hardware.
    \item Expert offload configurations (CPU offload for weights) introduce significant
    latency overhead that makes them unsuitable for interactive applications.
    \item 256K context performance is strong but not equivalent to 32K; tasks requiring
    full-context reasoning over 200K+ tokens should use Zen4 (1M) or wait for Zen4-Ultra-Long.
\end{itemize}

\section{Conclusion}

Zen4-Ultra demonstrates that the Zen MoDE framework scales reliably to one trillion total
parameters while maintaining practical inference through sparse activation of 96B parameters
per token. Targeted expert distillation produces genuinely specialized sub-networks that
outperform equivalent jointly-trained MoE experts, particularly on mathematics and scientific
reasoning. With new state-of-the-art results on MMLU (95.2\%), MATH (92.8\%), HumanEval
(96.1\%), GPQA Diamond (82.4\%), and MMMU (91.3\%), Zen4-Ultra represents the frontier
of what is achievable within practical deployment constraints in the Zen4 family.

\section*{Acknowledgments}

We thank the distributed systems teams who built and maintained the 16,384-H100 training
cluster, the annotation teams who produced the post-training datasets, and the evaluation
teams at Zoo Labs Foundation who designed and administered the safety evaluations.

\bibliographystyle{plain}
\bibliography{references}

\appendix

\section{Model Card}

\begin{table}[H]
\centering
\begin{tabular}{ll}
\toprule
\textbf{Field} & \textbf{Value} \\
\midrule
Model Name       & Zen4-Ultra \\
Version          & v2026.02 \\
Release Date     & February 2026 \\
Total Parameters & 1.02T \\
Active Parameters& 96B per token \\
Context Length   & 262,144 tokens \\
License          & Proprietary (API access) \\
API Endpoint     & \href{https://api.hanzo.ai/v1}{api.hanzo.ai/v1} \\
Documentation    & \href{https://papers.zenlm.org/zen4-ultra}{papers.zenlm.org/zen4-ultra} \\
Contact          & research@hanzo.ai \\
\bottomrule
\end{tabular}
\caption{Zen4-Ultra Model Card}
\end{table}

\end{document}
